% ============================================================
% EXEMPLES DE FIGURES ET IMAGES
% ============================================================

\subsection*{Figures et images}

\subsubsection*{Insérer une image simple}

Pour insérer une image, utilisez l'environnement \texttt{figure} :

\textbf{Code :}
\begin{verbatim}
\begin{figure}[H]
  \centering
  \includegraphics[width=0.5\textwidth]{chemin/image.png}
  \caption{Description de l'image}
  \label{fig:exemple1}
\end{figure}
\end{verbatim}

\textbf{Formats supportés :} PNG, JPG, PDF, SVG (avec conversion)

\subsubsection*{Contrôler la taille de l'image}

Différentes façons de spécifier la taille :

\begin{itemize}
	\item \texttt{width=0.5\textbackslash textwidth} : 50\% de la largeur du texte
	\item \texttt{width=8cm} : Largeur fixe de 8 cm
	\item \texttt{height=6cm} : Hauteur fixe de 6 cm
	\item \texttt{scale=0.75} : 75\% de la taille originale
	\item \texttt{width=0.8\textbackslash textwidth, height=6cm, keepaspectratio} : Respecte le ratio
\end{itemize}

\textbf{Exemple :}
\begin{verbatim}
\includegraphics[width=0.8\textwidth, keepaspectratio]{image.png}
\end{verbatim}

\subsubsection*{Plusieurs images côte à côte}

Pour placer plusieurs images sur la même ligne, utilisez \texttt{minipage} ou \texttt{subfigure} :

\textbf{Méthode 1 : minipage}
\begin{verbatim}
\begin{figure}[H]
  \centering
  \begin{minipage}{0.45\textwidth}
    \centering
    \includegraphics[width=\textwidth]{image1.png}
    \caption{Première image}
  \end{minipage}
  \hfill
  \begin{minipage}{0.45\textwidth}
    \centering
    \includegraphics[width=\textwidth]{image2.png}
    \caption{Deuxième image}
  \end{minipage}
\end{figure}
\end{verbatim}

\textbf{Méthode 2 : subfigure (avec package subcaption)}
\begin{verbatim}
\begin{figure}[H]
  \centering
  \begin{subfigure}{0.45\textwidth}
    \includegraphics[width=\textwidth]{image1.png}
    \caption{Sous-figure 1}
  \end{subfigure}
  \hfill
  \begin{subfigure}{0.45\textwidth}
    \includegraphics[width=\textwidth]{images/votre_image.png}
    \caption{Sous-figure 2}
  \end{subfigure}
  \caption{Comparaison de deux résultats}
  \label{fig:comparaison}
\end{figure}
\end{verbatim}

\subsubsection*{Graphiques TikZ}

Pour créer des diagrammes directement en LaTeX :

\textbf{Exemple simple :}
\begin{verbatim}
\begin{figure}[H]
  \centering
  \begin{tikzpicture}
    \draw[->] (0,0) -- (4,0) node[right] {$x$};
    \draw[->] (0,0) -- (0,3) node[above] {$y$};
    \draw[blue, thick] (0,0) -- (3,2);
    \fill[red] (3,2) circle (2pt);
  \end{tikzpicture}
  \caption{Graphique simple avec TikZ}
  \label{fig:tikz1}
\end{figure}
\end{verbatim}

\textbf{Diagramme de blocs :}
\begin{verbatim}
\begin{figure}[H]
  \centering
  \begin{tikzpicture}[node distance=2cm]
    \node[draw, rectangle] (input) {Entrée};
    \node[draw, rectangle, right of=input] (process) {Traitement};
    \node[draw, rectangle, right of=process] (output) {Sortie};
    \draw[->] (input) -- (process);
    \draw[->] (process) -- (output);
  \end{tikzpicture}
  \caption{Diagramme de flux}
  \label{fig:flux}
\end{figure}
\end{verbatim}

\subsubsection*{Référencer une figure}

Pour référencer une figure dans le texte :

\begin{verbatim}
Comme illustré à la Figure~\ref{fig:exemple1}, on observe...
La Figure~\ref{fig:comparaison} montre clairement que...
\end{verbatim}

\textbf{Note :} Le \texttt{\textasciitilde} crée une espace insécable entre "Figure" et le numéro.

\subsubsection*{Rotation d'image}

Pour faire pivoter une image :

\begin{verbatim}
\includegraphics[width=0.5\textwidth, angle=90]{image.png}
\end{verbatim}

\subsubsection*{Bonnes pratiques}

\begin{itemize}
	\item Toujours utiliser des chemins relatifs : \texttt{images/fig1.png}
	\item Préférer les formats vectoriels (PDF, SVG) pour les diagrammes
	\item Utiliser PNG pour les captures d'écran, JPG pour les photos
	\item Toujours ajouter une légende descriptive avec \texttt{\textbackslash caption}
	\item Toujours ajouter un label avec \texttt{\textbackslash label} pour pouvoir référencer
	\item Utiliser \texttt{[H]} du package \texttt{float} pour positionner exactement
\end{itemize}
